
\subsubsection{2.1 Calibration of Neural Networks and Language Models}

Guo et al. [2017] formalized ECE as the weighted average gap between predicted confidence and empirical accuracy across equal-width confidence bins, and demonstrated that modern neural networks are systematically overconfident. Temperature scaling---rescaling logits by a single learned scalar---reliably reduces ECE without affecting accuracy. ECE has since become the standard calibration metric in NLP and computer vision.

Minderer et al. [2021] extended calibration analysis to distribution shift in vision models, showing that ECE increases under standard image corruptions (ImageNet-C) and out-of-distribution benchmarks (ObjectNet). Their finding that distribution shift degrades calibration motivated the original hypothesis of this work. The results reported here complicate this picture: adversarial distribution shift in NLP does not uniformly degrade calibration, and the construction method of the adversarial benchmark determines whether calibration improves or worsens.

Kadavath et al. [2022] measured calibration of large language models using verbal probability elicitation on factual QA tasks, finding that RLHF-aligned models show better self-knowledge calibration than base models on clean benchmarks. Xiong et al. [2023] studied LLM confidence elicitation through verbal uncertainty expressions; Zhao et al. [2023] measured generation-based ECE for instruction-tuned LLMs. None of these works measure logit-based ECE on adversarial NLP benchmark splits.

\textbf{Research gap:} No prior work measures logit-based ECE on adversarial NLP benchmark splits for open-weight LLMs.

\subsubsection{2.2 Adversarial NLP Benchmarks}

Wang et al. [2021] introduced AdvGLUE, a multi-task adversarial benchmark constructed by applying human-designed perturbation methods (synonym substitution, word insertion, contextual perturbation) to GLUE examples to create inputs that preserve labels but cause model errors. AdvGLUE measures accuracy robustness only; calibration is not measured.

Nie et al. [2020] introduced ANLI (Adversarial NLI), constructed using a human-and-model-in-the-loop protocol: human annotators write NLI hypotheses that fool the current model, which is then retrained on the collected data for the next round. This produces three rounds of increasing difficulty (R1 < R2 < R3). ANLI is a model-in-the-loop benchmark---adversarial examples are selected to be difficult for the training model at each round---which makes it structurally different from AdvGLUE's human-only construction.

The distinction between human-adversarial (AdvGLUE) and model-in-the-loop adversarial (ANLI) construction is not typically highlighted in papers that use both benchmarks for accuracy evaluation. This work treats it as a primary independent variable, because this distinction is found to determine whether calibration degrades or improves.

\textbf{Research gap:} Adversarial benchmarks are used exclusively for accuracy evaluation; no prior work treats adversarial construction method as a variable in calibration analysis.

\subsubsection{2.3 RLHF Alignment and Calibration}

Reinforcement Learning from Human Feedback [Christiano et al., 2017; Ouyang et al., 2022] fine-tunes language models to align with human preferences. The effect of RLHF on calibration has been studied primarily on clean benchmarks. Kadavath et al. [2022] found RLHF models show better self-reported calibration; OpenAI's GPT-4 technical report \citep{OpenAI2023} notes calibration improvements from RLHF training. The conventional expectation is that RLHF improves calibration.

This work tests this expectation in the adversarial setting for the first time. The finding that RLHF moderation of adversarial calibration is benchmark-type conditional---helpful for ANLI, harmful for AdvGLUE---is not predicted by clean-benchmark RLHF calibration findings.

\textbf{Research gap:} RLHF's effect on calibration has been studied only on clean benchmarks; adversarial calibration moderation by RLHF is untested.

\subsubsection{2.4 Positioning}

This work occupies the intersection of calibration measurement, adversarial NLP benchmarks, and RLHF alignment effects. The conditional pattern documented here---where the same LLM can improve calibration on one adversarial benchmark while degrading it on another, depending on construction method---cannot be derived from any single prior thread.

